% ============================================
% Johns Hopkins Letter Template
% ============================================

\documentclass[11pt,letterpaper]{letter}

\usepackage{graphicx}
\usepackage{eso-pic}
\usepackage{geometry}
\usepackage{microtype}
\usepackage[utf8]{inputenc}
\usepackage[hidelinks]{hyperref}

\geometry{left=25mm,right=25mm,top=25mm,bottom=25mm}

\newcommand{\PutJHULogoFG}[1]{%
  \AddToShipoutPictureFG*{%
    \AtPageUpperLeft{%
      \hspace{10mm}%
      \raisebox{-38mm}{%
        \includegraphics[width=6cm]{#1}%
      }%
    }%
  }%
}

\signature{Decory Edwards}
\address{
Department of Economics \\
Johns Hopkins University \\
Baltimore, MD 21218 \\
\texttt{dedwar65@jh.edu}
}

\begin{document}
\PutJHULogoFG{Images/jhulogo.png}

\begin{letter}{
Search Committee \\
Department of Economics \\
Franklin and Marshall College \\
Lancaster, PA
}

\opening{Dear Members of the Search Committee,}

I am writing to apply for the Visiting Assistant Professor position in Economics at Franklin and Marshall College. I am a Ph.D.\ candidate in the Department of Economics at Johns Hopkins University (advisors: Christopher D.\ Carroll, Nicholas W.\ Papageorge, and Stelios Fourakis), and I expect to receive my degree in May 2026. My primary specializations are macroeconomic modeling, wealth and income dynamics, and computational economics. I have teaching experience at the university level and am committed to excellence in undergraduate education.

My research agenda focuses on the ability of macroeconomic models to generate empirically realistic distributions of wealth and income over the life cycle. A key driver of that work is modeling heterogeneity in the rate of return on assets, and understanding how return dispersion contributes to portfolio accumulation across households.

In my job market paper, I calibrate a life cycle heterogeneous agent model to match wealth moments from the Survey of Consumer Finances by allowing households to differ in their portfolio returns. The model helps explain observed wealth concentration and return dispersion. In a related research project using the Health and Retirement Study, I examine how trust influences both income growth and asset returns. I find a hump shaped relationship for trust and returns (consistent with the literature) and characterize the distribution of the persistent component of returns to net worth. These experiences have equipped me to frame research strategies and design modeling approaches, execute econometric and simulation analysis with large micro data sets, and translate results into clear and rigorous written deliverables for both academic and practitioner audiences.

I am particularly drawn to Franklin and Marshall’s liberal arts environment and its emphasis on foundational courses that introduce students to core economic perspectives and distributional questions. Courses such as ECO 103 and ECO 203 provide an opportunity to engage students with fundamental debates about value, inequality, and the allocation of resources. My research on wealth concentration, return heterogeneity, and public finance naturally connects to these themes. I strive to create an inclusive classroom environment that encourages analytical rigor while remaining accessible to students from diverse academic backgrounds. In a 3/3 teaching setting, I would be excited to contribute to introductory and intermediate courses and to offer electives that expose students to macroeconomics, public finance, and quantitative methods.

I am enthusiastic about relocating to Lancaster and fully engaging in the in person academic community at Franklin and Marshall. I value close interaction with undergraduates, mentorship, and the opportunity to contribute meaningfully to a department committed to inclusive excellence in teaching and student success. A visiting appointment would allow me to continue developing as both a teacher and scholar while contributing to the College’s strong tradition of liberal arts education.

I believe I would be a strong fit for this role because:
\begin{enumerate}
    \item My research on inequality and wealth dynamics aligns closely with the themes of value and distribution emphasized in the curriculum.
    \item I am committed to clear, inclusive, and engaging undergraduate teaching.
    \item I bring quantitative rigor and policy relevance that enrich both introductory and elective offerings.
\end{enumerate}

Thank you very much for your consideration. I would welcome the opportunity to contribute to the Department of Economics at Franklin and Marshall College.

\closing{Sincerely,}

\end{letter}
\end{document}
